%%%%%%%%%%%%%%%%%%%%%%%%%%%%%%%%%%%%%%%%%%%%%%%%%%%%%%%%%%%%%%%%%%%%%%%%%%%%%%%%
%%  DISCUSSION AND LIMITATIONS
%%%%%%%%%%%%%%%%%%%%%%%%%%%%%%%%%%%%%%%%%%%%%%%%%%%%%%%%%%%%%%%%%%%%%%%%%%%%%%%%
\chapter{Discussion and Limitations}\label{discussion}
The benchmark results show that the framework reliably translates natural-language instructions into coordinated actions. At the same time, the runs expose recurring failure modes and several open design questions. Together, these restrictions define the system’s usefulness today and set the agenda for the rest of this chapter.

% ------------------------------------------------------------------------------
\section{Perception and Depth Sensing}\label{sec:perception-depth}
The tuned \gls{sgbm} pipeline localizes objects under 1~m to within 5–10\% depth error, and this per-object depth estimate is robust. Taking object depth as the median disparity over the textured detection region averages out the zero-mean matching noise. That robustness does not carry over to the dense, per-pixel depth that \gls{vgn} requires. \gls{vgn}'s volumetric occupancy grid is built by back-projecting the full disparity map, and on the textureless surfaces that dominate the scene, per-pixel stereo matching fails systematically. Because that error is a systematic scale bias, median aggregation, which rescues object localization, cannot correct it. The result is sparse, discontinuous \gls{tsdf} volumes that mislead the grasp-quality network into predicting off-axis candidates or suppressing valid poses.

The heuristic geometric planner, therefore, acts as a conditional fallback, not the intended primary grasp source. It fires precisely when \gls{vgn} cannot return a usable pose, which is the common case here but reflects the degraded stereo input. The fallback itself is robust: it needs no dense 3D reconstruction, avoids the scale bias entirely, scores candidates on direction, contact geometry, and \gls{ik} quality, and always issues and closes a grasp.

Even so, the \gls{vgn} path is not abandoned. In simulation, it does not bypass the stereo sensor: its occupancy grid is built from the same \gls{sgbm} point cloud used for real images, and to work around the scale bias, the object position is taken from the simulator’s \texttt{WorldState} oracle rather than recovered from that cloud. A dedicated depth sensor behind the adapter’s existing \texttt{get\_depth\_frame} hook would supply that localization at the source and let an actual \gls{vgn} pose contribution be recovered, a deployment prerequisite we return to in the next chapter.

% ------------------------------------------------------------------------------
\section{Large Language Model Latency as a Control Bottleneck}
As Section~\ref{sec:communication} notes, end-to-end command latency breaks into three components, and \gls{vlm} parsing dominates. It is also the least predictable component: a short, unambiguous prompt parses in under 10~s, while a multi-robot task requiring complex coordination consumes the full thinking budget and approaches the 90-second request timeout.

This latency profile has two practical consequences. First, it precludes reactive control. The system cannot re-plan faster than the model generates tokens. A robot moving at 0.5~m/s travels 10~m during a median-length parse, so any obstacle or object displacement that occurs while the planner thinks is invisible to it. Replanning is therefore viable between motions, never during one. Second, planning cost and execution cost are decoupled. The negotiation protocol front-loads a single analysis–proposal–evaluation round of model calls per task, while leaving the per-operation execution rate unchanged. The longer wall-clock totals for the negotiation-enabled runs, therefore, reflect more completed work rather than slower planning and should not be read as a planning-speed effect.

% ------------------------------------------------------------------------------
\section{Inverse Kinematics Convergence at Workspace Boundaries}
The damped least-squares \gls{ik} solver converges reliably within the core of each robot’s reachable workspace but degrades near the boundary, where the Jacobian becomes ill-conditioned at near-singular configurations. The damping factor $\lambda = 0.2$ suppresses the most severe inversions, yet it cannot eliminate convergence failures for targets that fall exactly at the reachable radius or demand approach trajectories through kinematic singularities. Boundary targets thus remain intrinsically fragile even when the solver is tuned conservatively.

This degradation is an analytical property of the direct-\gls{ik} path, not something the benchmarks isolate, since the ablations toggle reflection and \gls{ros} movement, not the \gls{ik} solver alone. The B12 reflection ablation shows that transient motion failures at boundary coordinates occur and usually recover on a retry; because those moves were planned through \gls{ros}/MoveIt, they reflect transient planning and execution failures, not steady-state direct-\gls{ik} convergence. The practically important point is the non-determinism. Identical task sequences succeed in some runs and fail in others, an intermittent \textit{sometimes fails} error class that restricting command coordinates alone cannot guard against.

As B16 establishes, the advantage of \gls{ros}/MoveIt here is both capability and reliability. Its full-joint-space \gls{ompl} planning finds collision-free trajectories that direct \gls{ik} cannot, and its grasp routing avoids the direct path's side-approach failures (Section~\ref{sec:results-ros}). Extending this collision-aware planning from point-to-point motion to the full grasp pipeline is thus the clearest path to improving task robustness, since it directly reduces workspace-boundary and approach failures.

% ------------------------------------------------------------------------------
\section{Bimanual Manipulation and Synchronized Grasping}\label{sec:bimanual}
B7’s failure is easy to misdiagnose, so it is worth stating what the task asks for: a sequential two-arm grasp, each \texttt{grasp\_object} bracketed by a \texttt{signal}/\texttt{wait\_for\_signal} barrier, followed by a simultaneous lift to $y=0.15$~m. The two failure modes documented in Section~\ref{sec:results-dual}, the mis-realized barrier and the lift rejection, are therefore not missing operations. What the lift phase reveals is a lack of partner-aware target selection: each robot independently determines its end-effector target, so two valid lift poses can still collide, and the violations cluster tightly around 0.10–0.11~m.

A deeper, simulation-level constraint prevents genuine concurrent dual-grasping of a single object, regardless of target selection. The reparenting is itself a workaround. Friction-only grasping in Unity’s physics engine proved unreliable, with objects slipping or jittering out of the gripper, so on grasp \texttt{GripperController} rigidly attaches the object to the gripper transform via \texttt{SetParent} instead of relying on contact friction. A Unity~\cite{unity} \texttt{Transform} admits exactly one parent, however, so one object cannot attach to two grippers at once. The handoff path works only because it serializes. Before the second gripper attaches, it forces the first to release. This is also a reliability hazard: the receiver only re-parents once its attachment check succeeds (gripper within about 10~cm of the object), so if the source releases while the object is outside that window, it reverts to a friction-only hold and can fall onto the table. A true co-grasp would therefore require replacing the single-parent attachment model itself as a second blocker, independent of the missing partner-aware target selection.

Because these blockers are independent of the planner, B7 is partially unwinnable regardless of the plan’s quality. The run data still separates two outcomes worth keeping distinct. The planner reproduced the expected chain only approximately. Plan length varied from 7 to 10 operations, and 3 of the 5 \gls{magi} runs dropped the second synchronization barrier. The coordination safety layer, by contrast, worked as intended. Notably, the \texttt{SetParent} ceiling is never actually reached. Execution halts earlier at the handshake or the proximity check, so the attachment limit stays theoretical here. A benchmark that genuinely tests co-grasping would need either a dual-attachment representation or a configurable physics joint in place of the single reparenting transform, or it would score the lift based solely on target selection.

% ------------------------------------------------------------------------------
\section{Parse-Level Input Validation}\label{sec:parse-validation}
B9 shows that parser-level rejection of impossible instructions is unreliable for every model, as no backend rejects with any consistency except the weakest planner (\gls{gema42}), and it does so only by failing to emit a plan at all (Section~\ref{sec:results-robustness}).

This is not inherently unsafe. Because B9 stops at parse time (Section~\ref{sec:results-robustness}), no motion is dispatched, and at execution, the precondition and coordination guards would reject commands referencing an unknown robot or a missing object. But for models that do not reject at parse time, the safety guarantee against invalid inputs rests entirely on those downstream guards. A prompt that names a valid robot and object but combines them in an impossible way would likewise reach the execution layer and fail only at the grasp step, after wasting time on detection, motion, and contact verification.

A dedicated pre-parse validation step that checks robot identifiers against the registered robot set and object IDs against the live world state would catch the B9 class of errors deterministically before any model call and independently of model capability. The Robot Constitution already establishes the pattern of code-level pre-execution checks on AutoRT-generated tasks, but it validates kinematic constraints rather than identifier or object existence; extending it with those membership checks, applied to externally submitted commands, is the missing piece to close this benchmark limit.

% ------------------------------------------------------------------------------
\section{Knowledge Graph Context Injection}
B14 found that \gls{kg} context injection carries a small, consistent cost. It lowers parse-level success by about 1~pp and adds close to one flagged operation per run, on top of a disabled baseline that already carries 4–5 flagged operations per run. The graph, therefore, does not create hallucination where none existed; it slightly aggravates a residual floor that is already present and intrinsic to the parser/task combination.

The benchmark cannot localize this cost further, because it records a single combined count that sums two distinct checks, namely operations absent from the registry entirely, and \gls{kg}-aware operations whose parameters reference neither a \gls{kg} object id nor a \texttt{\$variable}. No intermediate condition separates the effect of the added prompt structure from that of the spatial facts’ semantic content, so the two cannot be told apart from this result alone.

Because that floor is independent of the graph, the targeted-injection redesign proposed in future work should be motivated by reduced prompt load rather than by the prospect of eliminating it.

% ------------------------------------------------------------------------------
\section{Simulation-to-Real Gap}\label{sec:sim-to-real}
The system is developed and validated entirely within Unity’s \texttt{ArticulationBody} simulation. Several properties of this environment do not transfer directly to a physical AR4 installation, as physical robot joints exhibit compliance, backlash, and friction that the \texttt{ArticulationBody} model does not reproduce, so the simulation-tuned \gls{pd} gains will need to be retuned against the real AR4’s inertia and damping for deployment.

Grasp verification compounds this deployment gap. In simulation, grasp force is read from the finger \texttt{ArticulationBody} joint forces and smoothed with a 5-frame moving average; the 5~N minimum-force threshold and 50~ms contact-duration requirement were tuned to reject simulation physics noise, not real fingertip normal forces, which depend on surface material, object deformability, and approach speed. Both must be revalidated against the real robot’s force response.

Object state is the larger substitution. Because the synthetic stereo overestimates depth by roughly 1.8$\times$ on Unity’s textureless surfaces, the grasp path reads object position, dimensions, and orientation from the simulator’s \texttt{WorldState} oracle instead of the point cloud (Section~\ref{sec:perception-depth}). \texttt{WorldState} is the interface a real perception stack would populate in its place. Detection already writes stereo-derived position back to it, so position closes with the depth-sensing upgrade; dimensions and orientation have no perception producer today and would need a segmentation-and-pose capability that the system does not yet build, not merely a better sensor.

The hardware adapter layer makes this deployment a configuration change. The adapters expose the same abstract interface as their Unity equivalents, so no operation code changes when switching modes. Motion is routed through the \gls{ros}/MoveIt backend, so trajectory execution on the physical AR4 depends on its \gls{ros} driver alone. Perception is already wired end-to-end. A capture bridge forwards the local camera adapter’s frames into the same image store that the pipeline reads in simulation. What remains is integration and validation on real hardware, namely depth sensing, force-torque feedback (currently used as a contact-duration heuristic), camera calibration, and validation on the real arms.

% ------------------------------------------------------------------------------
\section{Evaluation Scope and Coverage}
The B12 reflection benefit is real, and at twenty runs per condition, its latency cost is measurable. That cost falls entirely on the runs that fail: a disabled run that aborts at the 120-second per-step timeout averages 168~s, and enabling reflection roughly doubles a firing run's duration on top of that (Section~\ref{sec:results-reflection} breaks down the fired-versus-idle timings). This sharpens the trade-off: reflection reliably converts intermittent boundary-reach failures into successes ($+14$~pp task success) but is expensive precisely on the runs that need it most, where recovery latency is bounded.

Two coverage gaps remain. First, the evaluation cannot assess reflection over longer deployment windows or densely coupled task sequences. The five stress tasks were mutually independent, so an unrecovered failure cost only its own task, whereas in a tightly coupled sequence, a single failure would cascade downstream, and the mechanism’s value would be correspondingly larger. Second, B17 exercises only AutoRT’s own stages, the task generator, and the two-layer Robot Constitution, offline and in isolation. The full generate–filter–select loop, running end-to-end against a live, changing scene, remains unevaluated, and with it, selection-strategy behavior over long horizons, proposal diversity under repeated execution, and the constitution’s rejection profile as the scene evolves. Both regimes would first require defining success criteria for open-ended task generation, an open research question in itself.

Underlying all of these is the run budget. Most conditions were evaluated five times in a single simulated environment, with the B9, B11–B13, and \gls{magi} B16 sweeps extended to twenty; five runs is enough to expose deterministic failure modes but not to resolve small effects that lie within one \gls{sd}. A fuller evaluation would need more runs, harder tasks, hardware-in-the-loop testing, and injected faults to sharpen the latency takeaway.

